\chapter{Documentación HATEOAS}
\label{ch:hateoas}

\section{Framework elegido: Spring HATEOAS}

Se eligió \textbf{Spring HATEOAS} como framework para implementar el nivel de madurez
de Richardson correspondiente a hipermedia (HATEOAS), por su integración nativa con
Spring Boot: no requiere construir manualmente las URLs de los enlaces, sino que las
deriva directamente de los métodos del propio controlador mediante
\texttt{WebMvcLinkBuilder}, evitando que los enlaces queden desincronizados si cambia
una ruta.

\begin{lstlisting}[language=Java, caption={LibroController.java -- respuesta con enlace HATEOAS}]
@GetMapping("/{id}")
public ResponseEntity<EntityModel<LibroDTO>> obtenerPorId(@PathVariable Long id) {
    Libro libro = libroService.obtenerPorId(id);
    LibroDTO dto = libroMapper.toDTO(libro);

    EntityModel<LibroDTO> model = EntityModel.of(dto);
    model.add(linkTo(methodOn(LibroController.class)
        .obtenerPorId(id)).withSelfRel());
    model.add(linkTo(methodOn(LibroController.class)
        .listar()).withRel("listar"));

    return ResponseEntity.ok(model);
}
\end{lstlisting}

Cada respuesta de un recurso individual incluye al menos un enlace \textit{self}, y en
varios casos un enlace adicional hacia el recurso relacionado más natural (por ejemplo,
de un libro hacia el listado completo, o de una reseña hacia el promedio de
calificaciones del libro correspondiente).

\section{Navegabilidad de la API}

El siguiente ejemplo, tomado de una respuesta real capturada durante la verificación en
Docker, muestra la forma exacta en que Spring HATEOAS serializa un recurso individual
(formato HAL): los campos del recurso se aplanan al nivel superior del JSON, junto al
objeto \texttt{\_links}:

\begin{lstlisting}[language=json, caption={Respuesta real de \texttt{POST /api/v1/catalogo}}]
{
  "_links": {
    "self": { "href": "http://localhost:8081/api/v1/catalogo/3" }
  },
  "id": 3,
  "titulo": "Cien Anos de Soledad",
  "autor": "Gabriel Garcia Marquez",
  "precio": 15990.0
}
\end{lstlisting}

Para una colección, en cambio, el formato HAL cambia de estructura: en lugar de un
arreglo plano, la respuesta envuelve los elementos dentro de \texttt{\_embedded}:

\begin{lstlisting}[language=json, caption={Respuesta real de un listado con HATEOAS}]
{
  "_embedded": {
    "historialEstadoList": [
      { "id": 1, "estado": "PREPARACION", "envioId": 16 },
      { "id": 2, "estado": "EN_CAMINO",  "envioId": 16 }
    ]
  },
  "_links": {
    "self": { "href": "http://localhost:8088/api/v1/envios/16/historial" }
  }
}
\end{lstlisting}

Esta diferencia estructural entre un recurso individual (aplanado) y una colección
(\texttt{\_embedded}) no es cosmética: un consumidor de la API que espere un arreglo
plano de JSON (como un \texttt{List<X>} tipado en un cliente Feign) fallará al
deserializar una respuesta HATEOAS de colección. Este problema se presentó realmente al
extender HATEOAS al microservicio de inventario: el \texttt{InventarioClient} que
consume \texttt{pedidos} esperaba \texttt{List<InventarioRentDTO>}, y tuvo que
actualizarse a \texttt{CollectionModel<EntityModel<InventarioRentDTO>>} para reflejar
la forma real de la respuesta:

\begin{lstlisting}[language=Java, caption={InventarioClient.java (pedidos) -- adaptado a la forma HATEOAS real}]
// inventario expone esta lista envuelta en HATEOAS (CollectionModel), no como
// un arreglo plano -- hay que declarar el tipo tal cual lo devuelve el productor.
@GetMapping("/api/v1/inventario/libro/{libroId}")
CollectionModel<EntityModel<InventarioRentDTO>> obtenerStockPorLibro(
    @PathVariable("libroId") Long libroId);
\end{lstlisting}

Este detalle refuerza por qué la navegabilidad HATEOAS debe considerarse desde el
diseño del contrato entre microservicios, y no agregarse como una capa cosmética al
final del desarrollo.
